% Part: normal-modal-logic
% Chapter: completeness
% Section: lindenbaums-lemma

\documentclass[../../../include/open-logic-section]{subfiles}

\begin{document}

\olfileid[fr]{nml}{com}{lin}

\olsection{Lemme de Lindenbaum}

Le lemme de Lindenbaum établit que tout ensemble $\Sigma$-cohérent de
!!{formula}s est contenu dans au moins un ensemble $\Sigma$-cohérent
\emph{complet}. La construction du modèle canonique montrera qu'à
chaque ensemble $\Sigma$-cohérent complet~$\Delta$ correspond un monde
du modèle canonique où les !!{formula}s vraies sont exactement celles
de~$\Delta$. Le lemme de Lindenbaum garantit donc que, pour tout
ensemble $\Sigma$-cohérent, il existe un monde du modèle canonique où
toutes les formules de cet ensemble sont vraies.

\begin{thm}[Lemme de Lindenbaum]\ollabel{thm:lindenbaum}
  Si $\Gamma$ est $\Sigma$-cohérent, il existe un ensemble
  $\Sigma$-cohérent complet $\Delta$ qui prolonge~$\Gamma$.
\end{thm}

\begin{proof}
Énumérons toutes les formules du langage sous la forme $!A_0$,
$!A_1$, \dots{} ; les répétitions sont permises. On peut, par exemple,
commencer par $\Obj p_0$ et, à chaque étape $n \ge 1$, énumérer les
formules, en nombre fini, de longueur au plus~$n$ qui n'utilisent que
les variables parmi $\Obj p_0$, \dots,~$\Obj p_n$. Nous définissons
des ensembles de !!{formula}s $\Delta_n$ par récurrence sur~$n$, puis
posons $\Delta = \bigcup_n \Delta_n$. D'abord, $\Delta_0 = \Gamma$.
Une fois $\Delta_n$ défini, nous définissons $\Delta_{n+1}$ par :
\[
\Delta_{n+1} =
\begin{cases}
  \Delta_n \cup \{!A_n\}, & \text{si $\Delta_n \cup \{ !A_n\}$
    est $\Sigma$-cohérent ;} \\
  \Delta_n \cup \{ \lnot !A_n\}, & \text{sinon.}
\end{cases}
\]
Posons alors $\Delta = \bigcup_{n=0}^\infty \Delta_n$.

Il faut vérifier que cette définition produit bien un ensemble $\Delta$
ayant les propriétés requises, à savoir $\Gamma \subseteq \Delta$ et
que $\Delta$ est $\Sigma$-cohérent complet.

L'inclusion $\Gamma \subseteq \Delta$ est immédiate : par
construction, $\Delta_0 \subseteq \Delta$ et $\Delta_0 = \Gamma$. En
fait, $\Delta_n \subseteq \Delta$ pour tout~$n$, puisque $\Delta$ est la
réunion de tous les~$\Delta_n$. (À chaque étape, nous ajoutons une
!!{formula} à l'ensemble déjà construit, donc $\Delta_n \subseteq
\Delta_{n+1}$. Par transitivité de $\subseteq$, on a aussi $\Delta_n
\subseteq \Delta_{m}$ dès que $n \le m$.) À chaque étape de la
construction, nous ajoutons soit $!A_n$, soit $\lnot !A_n$, et toute
!!{formula} figure au moins une fois dans la liste des~$!A_n$. Ainsi,
pour tout $!A$, soit $!A \in \Delta$, soit $\lnot !A \in \Delta$ ; par
définition, $\Delta$ est donc complet.

Enfin, il reste à montrer que $\Delta$ est $\Sigma$-cohérent. Nous
établissons pour cela (a) que si $\Delta$ était $\Sigma$-incohérent,
alors un certain $\Delta_n$ serait $\Sigma$-incohérent, et (b) que tous les
$\Delta_n$ sont $\Sigma$-cohérents.

Supposons donc que $\Delta$ soit $\Sigma$-incohérent. On a alors
$\Delta \Proves[\Sigma] \lfalse$ : il existe des formules $!A_1$,
\dots,~$!A_k \in \Delta$ telles que $\Sigma \Proves !A_1 \lif (!A_2
\lif \cdots (!A_k \lif \lfalse)\dots)$. Puisque $\Delta =
\bigcup_{n=0}^\infty \Delta_n$, chaque $!A_i \in \Delta_{n_i}$ pour un
certain~$n_i$. Prenons $n=\max(\{0\}\cup\{n_i : 1\le i\le k\})$. Comme $n_i \le
n$, on a $\Delta_{n_i} \subseteq \Delta_n$. Toutes les $!A_i$
appartiennent donc à un même~$\Delta_n$. Il s'ensuit que $\Delta_n
\Proves[\Sigma] \lfalse$, autrement dit que $\Delta_n$ est
$\Sigma$-incohérent.

Pour montrer que chaque $\Delta_n$ est $\Sigma$-cohérent, procédons
par récurrence sur~$n$. On a $\Delta_0 = \Gamma$, et $\Gamma$ est
$\Sigma$-cohérent par hypothèse ; la propriété vaut donc pour $n = 0$.
Supposons-la vraie au rang $n$, c'est-à-dire supposons $\Delta_n$
$\Sigma$-cohérent. L'ensemble $\Delta_{n+1}$ est soit
$\Delta_n \cup \{!A_n\}$ si celui-ci est $\Sigma$-cohérent, soit
$\Delta_n \cup \{\lnot!A_n\}$ sinon. Dans le premier cas,
$\Delta_{n+1}$ est évidemment $\Sigma$-cohérent. Or, d'après la
\olref[prf][con]{prop:consistencyfacts}\olref[prf][con]{prop:consistencyfacts-c},
l'un au moins des ensembles $\Delta_n \cup \{!A_n\}$ et
$\Delta_n \cup \{\lnot!A_n\}$ est cohérent ; $\Delta_{n+1}$ l'est
donc dans l'autre cas également.
\end{proof}

\begin{cor}\ollabel{cor:provability-characterization}
  $\Gamma \Proves[\Sigma] !A$ si et seulement si $!A \in \Delta$ pour
  tout ensemble $\Sigma$-cohérent complet $\Delta$ qui prolonge
  $\Gamma$ (y compris lorsque $\Gamma = \emptyset$, ce qui donne une
  autre caractérisation du système modal $\Sigma$).
\end{cor}

\begin{proof}
  Supposons $\Gamma \Proves[\Sigma] !A$ et soit $\Delta$ un ensemble
  $\Sigma$-cohérent complet quelconque qui prolonge $\Gamma$. Si $!A
  \notin \Delta$, alors, par maximalité, $\lnot!A \in \Delta$. On a
  donc $\Delta \Proves[\Sigma] !A$ par monotonie et $\Delta
  \Proves[\Sigma] \lnot!A$ par réflexivité, ce qui rendrait $\Delta$
  incohérent. Réciproquement, si $\Gamma \Proves/[\Sigma] !A$, alors
  $\Gamma \cup \{ \lnot!A\}$ est $\Sigma$-cohérent ; par le lemme de
  Lindenbaum, il existe un ensemble cohérent complet $\Delta$ qui
  prolonge $\Gamma \cup \{ \lnot!A \}$. Par cohérence, $!A
  \notin \Delta$.
\end{proof}
\end{document}
